\section{授权历史修复}
\label{sec:history}

\subsection{在原始坐标处读取资金引用}
\label{sec:protocol-reader}

归档操作者重新检查一笔已执行支付，以说明缺失输入如何获得资金、发行方是否已获偿还。对于 $T_B$，记录应标明从 $T_A$ 接受的输出 $x$、针对 $G_A$ 的赔付扣款（若有），以及当前义务状态。操作者可以导出原始正文和授权修订正文，后者在同一输入位置记录垫付。

该服务已验证并执行前缀 $H$。它在同一 $V_H$ 中获取授权正文、精确的 Revision 与 Task、永久决定及当前义务状态；使用授权开口重建区块和 part，对照保存的原始提交证据检查完整 BlockID，并连同决定和修订上下文导出原始与修订正文。提交证据检查与既有区块身份的相容性；已执行命令与精确 Task 授权替换字节。投票同时覆盖头部哈希与 part-set header，因此只保持头部哈希并不足够。如果 Alice 的来源随后偿还发行方，扣款引用仍记录历史垫付，而义务状态变为 Recovered。

因此，C3 的正确性目标不同于 C2 所结算的金额：即使备付扣款正确，也不得授权一个将该输入归因于另一发行方扣款的正文。修复必须同时保持被允许的资金引用，以及固定的付款与区块身份；依据决定进行的检查拒绝此类候选。原型实现了这一使用 Next 接口的本地工作流（表~\ref{tab:c3-function}）；它依赖已执行前缀，而不将导出的正文视为独立的远程证明。

位于 $(h,j,k)$ 的读取返回 $(H,h,j,k,F,r,d,s_H)$，由现有函数在 $V_H$ 上的一个读事务内组装而成。其中 $F$ 是资金表示，$r$ 是其修订，$d$ 是决定或缺失标记，$s_H$ 是义务状态。该逻辑元组描述本地组合，而不是已部署的网络接口。索引 $j$ 对原始交易列表计数，不会因过滤维护命令或失败交易而重新编号。逻辑区块体来自 $V_H$ 中已提交的修订，仅当不存在修订时才回退到原始区块，因此物理安装进度，或重放期间已持有较晚版本的区块存储，都不改变答案。存在但无法解码的修订是错误，绝不回退到原始区块。缺失的历史被报告为不可用，而不被解读为不存在赔付；该契约假设相关历史未被裁剪。在完整、已验证、未裁剪的 $V_H$ 中，缺失决定意味着截至 $H$ 不存在决定。



因此，读取者将永久决定与义务状态连同资金表示一并查阅。待表示条目是可移除的工作状态，而不是赔付已发生的持久证据。这些表示或安装变更都不会向消费方计入第二笔本金。

第~\ref{sec:model-records}~节介绍的状态维度保持相互分离。

索引或普通物化视图保持原始正文，同样能回答经济问题。经过另行认证的修订副本也可保留旧 ID 查询，Reparo 展示了这种访问模式~\cite{reparo}。C3 则使授权修订正文自身仍满足原完整 BlockID。这是其记录层服务，而非支付最终性的额外条件。表示服务停止时，经济闭合仍然推进；当前构建依然使用 CH 交易与 part 编码。

\subsection{表示命令与批次}
\label{sec:protocol-batch}

一个 \texttt{RepairBatch} 针对一个历史区块，按规范顺序含 1--32 个条目，每个条目指明一个输出、一个交易索引、一个输入索引和一个新的输入开口；输出与槽位唯一。该批次固定基础修订以及先前与替换后的规范区块体摘要，携带 part 开口，并将逻辑修订从 $\mathsf{Base}$ 推进到 $\mathsf{Base}+1$。单个 \texttt{RepairInput} 是其单条目形式，并携带被替换的交易字节。两种命令都只接受具有已提交赔付决定的输出，从决定记录中获取金额、发行方和替换引用，并且不写入会计状态。

公开执行对照其当前状态重复所有检查。基础已变化、某条目没有决定，或某槽位不再持有 \texttt{OriginalFunding}，都会拒绝整个批次，且确切的成功重放在基础验证之前即被识别。由于每条命令都要求其槽位中为 \texttt{OriginalFunding}，不同的批次或单条目入口不能对同一槽位替换两次，命令与修订标识保持特定于所选的打包方式。经认证的结果指明该命令及其是否被应用；跟随者将其映射为空的经济效果，而赔付的经济效果仅来自决定结果。偿还属于来源交易自身的执行结果。

\subsection{保持身份的安装与重放}
\label{sec:protocol-installation}

四种身份被区分开。支付身份 $\mathsf{id}(T)$ 固定授权与承诺，且所有者签名绑定它；它不是共识交易哈希。格式良好的 \texttt{UTXO4CH} 信封的共识 $\mathit{Tx.Hash}$ 是对其头部与固定字段的 SHA-256，不检查任何支付有效性。$\mathit{DataHash}$ 是共识交易哈希之上的 Merkle 根，而包含性证明不授予修订权限。完整 BlockID $(H_b,P_b)$ 保留区块哈希与原始 part-set 头，执行与安装都比较这两者。Part 承诺是绑定到链、密钥、高度和 part 索引的变色龙承诺，被更改的 part 携带指向其原始承诺的有效开口。原始共识 part 携带 revision~0 与 opening~1；原始验证要求正确的高度与这些值，且应用正常验证原始支付。被修订的 part 不能冒充新的原始执行，只有已提交的 Task 才授权改写。重放使用保留的原始命令，且从不将被修订的区块体作为新支付执行。

公开的修订与 Task 授权一个确切的新区块体及其 part 证据，安装器在获取物理修订锁之前，从一个已提交的读视图中复制它们。对于目标修订 $r$，较低的物理版本被验证并安装，相等的版本要求材料相同，较高的版本则不经回滚而覆盖旧任务。原始区块被保留，part 与本地版本元数据被原子地写入。安装可以跳过中间的物理版本，并保留每条经济命令。队列进度保持全局且连续：对于条目 $A_1,B_1,A_2$，在处理 $A_1$ 时安装 $A_2$ 会使 $B_1$ 保持未完成。逻辑执行读取已提交的 Canonical 修订，因此每种安装调度都留下相同的经济历史。当某个修订的区块体已在某副本上被物理安装，该修订即在该副本处被\emph{物化}。

重放按已提交的顺序执行保留的原始区块，包括支付、赔付决定和表示命令，并重建该序列的应用哈希。既有的后继交易保持其授权、输出引用和区块身份。仅赔付决定就以仅追加的公开记录在经济上关闭义务。表示增加了一个接口：在消费方的原始坐标 $(h,j,k)$ 处，已存储的区块指明指定输入的备付扣减，并仍能对照原始 BlockID 验证。读取者仍读取永久决定与义务状态以获得经济状态，被修订的字段并不取代它们。批处理在一个区块内共享表示工作，直接安装合并历史写入，二者都不为正常认证增加一轮交互。第~\ref{sec:eval-reader}~节报告了该接口的功能区别，以及相对于优化后决定索引的读取、存储与维护权衡。

\noindent\textbf{密码学实例化。}
该实现使用 RSA 变色龙哈希，模数为 2048 位，$e=65537$，并使用带上下文绑定的 SHAKE256 全域哈希映射到 $\mathbb Z_N^*$。一个可信 dealer 使用 CIRCL~v1.6.5~\cite{circl} 生成三选四的门限 RSA 密钥，遵循 Shoup 的门限构造~\cite{shoup}。公开的适配比值不传递编辑权限；传递编辑权限的是已提交的决定、确切的 Task 以及完整身份守卫。补充材料规定了编码、参数假设和原始群适配接口。
